%% sources: 05 discussions + 06 key findings/RQ -> target ~2 pages

\section{Discussion}\label{sec:discussion}

\paragraph*{Co-optimization is meaningful but uneven.}
Jointly optimizing style and tiles improves transfer size and rendering substantially, but unevenly.
Verbose institutional styles with hundreds of layers benefit far more than compact basemaps, and overview scenarios more than street-level panning.
The reason is structural.
Tile shaving removes only what a style never references, so its headroom is limited by how much of the source schema the style ignores.
This lets lean styles such as Toner shave the most and broad ones such as the two Bright variants the least (\cref{sec:results:data}).
There is thus no single optimal tile.
The optimal encoding is a function of the rendering style, so the same source data yields different ideal tiles for different consumers.

\paragraph*{Style-level and data-level passes compose additively.}
Against our prior expectations, the style-level passes (expression and structural) and the data-level passes (shaving and encoding) compose additively.
We expected them to compound, reasoning that style-level dead elimination would narrow the references feeding the tile-shaving advisory.
The data refutes this.
Shaving driven by the optimized style removed essentially the same bytes as shaving driven by the unoptimized one, the style passes adding at most half a percent (\cref{sec:results:combined}).
The advisory is based on which source-layers and columns a style references, not by how its expressions are written, so simplification leaves the shave target untouched, and on data volume and load time the two groups are simply additive.
Within the style level the picture is synergistic, since expression-level contradiction detection feeds structural dead elimination, which in turn removes the gaps between otherwise-mergeable layers.

\paragraph*{Which level dominates depends on the bottleneck.}
Because the three levels target different costs, dominance depends on the bottleneck.
The data-level passes dominate transfer size and load time, since tiles are fetched throughout interaction whereas the style is parsed once, so on a constrained mobile network they are what matters.
The structural passes dominate sustained frame rate, since once tiles are resident on the GPU the per-frame cost is draw-call count and per-feature expression evaluation rather than data size.
The expression passes shrink the style document and its parse time but rarely move either end-to-end metric on their own.
The axes meet only at overview zoom, where shaving entire source-layers reduces the features rendered per frame.

\paragraph*{Implications for map-serving infrastructure.}
Style-dependent tiles also fragment edge caches.
Two styles that once shared identical \ac{MVT} tiles now fetch distinct shaved variants, so a \ac{CDN} must key its cache on (tile, style) pairs.
Shaving a style pays off when its per-tile rendering saving exceeds the origin-fetch cost weighted by the cache hit rate the style forfeits by leaving the shared pool, so the decision turns on a style's share of pooled traffic, not on the number of styles served.
A dominant style already warms most of the pool, so any positive per-tile saving justifies shaving it.
A unpopular style cannot keep a private cache warm and reverts to cold-start behavior.
The optimal deployment is thus hybrid, shaving the few canonical styles and leaving the long tail on raw \acp{MVT}.

\paragraph*{Generalization to unseen styles.}
The cross-style evaluation argues these gains are not artifacts of our corpus.
On twelve previously unseen styles under static passes alone, the optimizer produced a meaningful reduction for every style and regressed none, the expression passes firing on virtually every input, as expected of any authored style.
We found no significant correlation between a style's initial complexity and its achievable reduction at this sample size, so the gains do not depend on a style being unusually verbose.
We read these figures as a lower bound, since the stats-driven and data-level passes need tile access the generalization corpus lacks.

\paragraph*{Limitations.}
The measurements carry the usual limits of a single-configuration benchmark.
All rendering ran on one hardware and browser configuration with ANGLE/Vulkan, so results may differ on other hardware or graphics backends, and the 15-run median may not capture rare events such as thermal throttling.
Network conditions were simulated through a fixed bandwidth throttle and a caching proxy rather than real \ac{CDN} latency, which also masks real-world cache behavior.
The ablation design means each pass's reported marginal contribution depends on the order of preceding passes, and tail metrics (p99 frame time, jank count) carry higher variance than central-tendency metrics.
Because the \num{18} scenarios share hardware, tile data, and styles, they are not fully independent, so we read the reported $p$-values as descriptive evidence of large effect sizes rather than strict inferential tests, and apply no multiple-comparisons correction across the full ablation grid.
Visual-equivalence checks across the corpus, all scenarios, and zoom levels 0 to 14 confirmed every pass stays within the per-test pixel-difference tolerance of \cref{sec:methods:problem}.

The conclusions are further limited in four ways.
First, the per-pass contributions are established on open-source styles only.
Three styles were excluded from end-to-end rendering by an upstream bug in MapLibre's style-migration tooling regarding mixed legacy and non-legacy expressions, and proprietary styles such as Mapbox Streets\footnote{\url{https://www.mapbox.com/maps/streets}, last accessed 2026-09-14} and Esri's vector basemaps\footnote{\url{https://developers.arcgis.com/documentation/mapping-and-location-services/mapping/basemap-layers/}, last accessed 2026-09-14} could not be evaluated because their terms of service forbid it.
Second, all tile-level results use MapLibre GL JS and the \ac{OMT} schema.
Schemas with other property distributions may yield different shaving gains, and the 10 benchmark styles, though diverse, may not cover all authoring patterns.
Third, geometry simplification was excluded from the optimizer because lossy vertex removal would violate the visual-equivalence criterion of \cref{def:visual_equivalence}.
The reported reductions therefore understate what a production pipeline could reach by composing the optimizer with an upstream simplifier.
Fourth, the jank-count metric worsened across the ablation only at sustained frame rates above \num{2000} \ac{FPS}, where garbage-collection pauses of fixed duration register as proportionally large frame-time outliers.
Consumer hardware caps frame rates well below this regime, so we do not read the worsening as a rendering regression.
